% GENERATED FILE. Edit canonical lessons, then run npm run generate:books.
% canonical-source-hash: fnv1a64:f16c56c0a6c39bbc
% canonical-lessons: MW-C110-batno, MW-C110-gamavno, MW-C110-pugno, MW-C110-urno, MW-C110-tolno

\chapter{To share out, To mislay, To reach, To fly, To weigh}
\label{ch:mw-to-share-out-to-mislay-to-reach-to-fly-to-weigh}
\hlchaptermodality{110}

\begin{chapteropening}
\textbf{By the end of this chapter:} \emph{I can say to share out, to mislay, to reach, to fly and to weigh.}

Complete the last lesson of chapter 110: I can say to share out, to mislay, to reach, to fly and to weigh.
\end{chapteropening}

\section[\texorpdfstring{\mw{बांटणो} (bā̃ṭṇo)}{बांटणो (bā̃ṭṇo)}]{\mw{बांटणो} (bā̃ṭṇo) — to share out}

\label{lesson:MW-C110-batno}



\begin{quote}
Before the new one: say the Marwadi for to lose (a game), then the Marwadi for to hit.
\end{quote}

\subsection*{You'll want to know: \mw{बांटणो}}
\textbf{\mw{बांटणो}} — \emph{bā̃ṭṇo} — \textquotedblleft{}to share out\textquotedblright{}.

\begin{cousinweb}
\textbf{\mw{बांटणो}} (\emph{bā̃ṭṇo}) is a tadbhava, a word inherited from Sanskrit and reshaped by sound change. It comes from \emph{vaṇṭati}, \textquotedblleft{}shares\textquotedblright{}. CDIAL lists Old Marwari \emph{bāṃṭaï} under it (CDIAL 11238).
\end{cousinweb}

\begin{grammarlens}[title={What you've built}]
One more word for this chapter.
\end{grammarlens}

\subsection*{Your turn}
\begin{itemize}
\raggedright
  \item \emph{Say it:} \emph{bā̃ṭṇo}
  \item \emph{Say it:} \emph{bā̃ṭṇo}, once more
  \item \emph{Say it:} say \emph{bā̃ṭṇo}
  \item \emph{Recall:} say the Marwadi for to lose (a game), then the Marwadi for to hit, then say \emph{bā̃ṭṇo} again
\end{itemize}

\begin{tcolorbox}[breakable,colback=teal!4,colframe=teal!35!black,title={Before you move on}]
What does \mw{बांटणो} mean? (\textquotedblleft{}To share out\textquotedblright{}.) Say it once more.
\end{tcolorbox}
\section[\texorpdfstring{\mw{गमावणो} (gamāvṇo)}{गमावणो (gamāvṇo)}]{\mw{गमावणो} (gamāvṇo) — to mislay, to lose (a thing)}

\label{lesson:MW-C110-gamavno}



\begin{quote}
Before the new one: say the Marwadi for to hit, then the Marwadi for to share out.
\end{quote}

\subsection*{You'll want to know: \mw{गमावणो}}
\textbf{\mw{गमावणो}} — \emph{gamāvṇo} — \textquotedblleft{}to mislay, to lose (a thing)\textquotedblright{}.

\begin{cousinweb}
\textbf{\mw{गमावणो}} (\emph{gamāvṇo}) is a tadbhava, a word inherited from Sanskrit and reshaped by sound change. It comes from \emph{gamayati}, \textquotedblleft{}causes to go\textquotedblright{}. CDIAL lists Old Marwari \emph{gavī} under it (CDIAL 4028). The Old Marwari row means \textquotedblleft{}lost\textquotedblright{}; the sense moved from \textquotedblleft{}send away\textquotedblright{} to \textquotedblleft{}lose\textquotedblright{}.
\end{cousinweb}

\begin{grammarlens}[title={What you've built}]
One more word for this chapter.
\end{grammarlens}

\subsection*{Your turn}
\begin{itemize}
\raggedright
  \item \emph{Say it:} \emph{gamāvṇo}
  \item \emph{Say it:} \emph{gamāvṇo}, once more
  \item \emph{Say it:} say \emph{gamāvṇo}
  \item \emph{Recall:} say the Marwadi for to hit, then the Marwadi for to share out, then say \emph{gamāvṇo} again
\end{itemize}

\begin{tcolorbox}[breakable,colback=teal!4,colframe=teal!35!black,title={Before you move on}]
What does \mw{गमावणो} mean? (\textquotedblleft{}To mislay, to lose (a thing)\textquotedblright{}.) Say it once more.
\end{tcolorbox}
\section[\texorpdfstring{\mw{पूगणो} (pūgṇo)}{पूगणो (pūgṇo)}]{\mw{पूगणो} (pūgṇo) — to reach, to arrive}

\label{lesson:MW-C110-pugno}



\begin{quote}
Before the new one: say the Marwadi for to share out, then the Marwadi for to mislay.
\end{quote}

\subsection*{You'll want to know: \mw{पूगणो}}
\textbf{\mw{पूगणो}} — \emph{pūgṇo} — \textquotedblleft{}to reach, to arrive\textquotedblright{}.

\begin{cousinweb}
\textbf{\mw{पूगणो}} (\emph{pūgṇo}) is a tadbhava, a word inherited from Sanskrit and reshaped by sound change. It comes from \emph{pūryate}, \textquotedblleft{}is filled\textquotedblright{}. CDIAL gives no Marwari form in this entry (CDIAL 8342); the match is through Gujarati \emph{pugvũ}. The Gujarati row means \textquotedblleft{}to reach\textquotedblright{}, as \textbf{\mw{पूगणो}} (\emph{pūgṇo}) does; Old Gujarati \emph{pūjaï} means \textquotedblleft{}reaches\textquotedblright{}.
\end{cousinweb}

\begin{grammarlens}[title={What you've built}]
One more word for this chapter.
\end{grammarlens}

\subsection*{Your turn}
\begin{itemize}
\raggedright
  \item \emph{Say it:} \emph{pūgṇo}
  \item \emph{Say it:} \emph{pūgṇo}, once more
  \item \emph{Say it:} say \emph{pūgṇo}
  \item \emph{Recall:} say the Marwadi for to share out, then the Marwadi for to mislay, then say \emph{pūgṇo} again
\end{itemize}

\begin{tcolorbox}[breakable,colback=teal!4,colframe=teal!35!black,title={Before you move on}]
What does \mw{पूगणो} mean? (\textquotedblleft{}To reach, to arrive\textquotedblright{}.) Say it once more.
\end{tcolorbox}
\section[\texorpdfstring{\mw{उडणो} (uṛṇo)}{उडणो (uṛṇo)}]{\mw{उडणो} (uṛṇo) — to fly}

\label{lesson:MW-C110-urno}



\begin{quote}
Before the new one: say the Marwadi for to mislay, then the Marwadi for to reach.
\end{quote}

\subsection*{You'll want to know: \mw{उडणो}}
\textbf{\mw{उडणो}} — \emph{uṛṇo} — \textquotedblleft{}to fly\textquotedblright{}.

\begin{cousinweb}
\textbf{\mw{उडणो}} (\emph{uṛṇo}) is a tadbhava, a word inherited from Sanskrit and reshaped by sound change. It comes from \emph{uḍḍīyate}, \textquotedblleft{}flies up\textquotedblright{}. CDIAL lists Old Marwari \emph{ūḍaï} under it (CDIAL 1697).
\end{cousinweb}

\begin{grammarlens}[title={What you've built}]
One more word for this chapter.
\end{grammarlens}

\subsection*{Your turn}
\begin{itemize}
\raggedright
  \item \emph{Say it:} \emph{uṛṇo}
  \item \emph{Say it:} \emph{uṛṇo}, once more
  \item \emph{Say it:} say \emph{uṛṇo}
  \item \emph{Recall:} say the Marwadi for to mislay, then the Marwadi for to reach, then say \emph{uṛṇo} again
\end{itemize}

\begin{tcolorbox}[breakable,colback=teal!4,colframe=teal!35!black,title={Before you move on}]
What does \mw{उडणो} mean? (\textquotedblleft{}To fly\textquotedblright{}.) Say it once more.
\end{tcolorbox}
\section[\texorpdfstring{\mw{तोलणो} (tolṇo)}{तोलणो (tolṇo)}]{\mw{तोलणो} (tolṇo) — to weigh}

\label{lesson:MW-C110-tolno}



\begin{quote}
Before the new one: say the Marwadi for to reach, then the Marwadi for to fly.
\end{quote}

\subsection*{You'll want to know: \mw{तोलणो}}
\textbf{\mw{तोलणो}} — \emph{tolṇo} — \textquotedblleft{}to weigh\textquotedblright{}.

\begin{cousinweb}
\textbf{\mw{तोलणो}} (\emph{tolṇo}) is a tadbhava, a word inherited from Sanskrit and reshaped by sound change. It comes from \emph{tolayati}, \textquotedblleft{}lifts, weighs\textquotedblright{}. CDIAL gives no Marwari form in this entry (CDIAL 5979); the match is through Hindi \emph{tolnā} and Gujarati \emph{toḷvũ}.
\end{cousinweb}

\begin{grammarlens}[title={What you've built}]
One more word for this chapter.
\end{grammarlens}

\subsection*{Your turn}
\begin{itemize}
\raggedright
  \item \emph{Say it:} \emph{tolṇo}
  \item \emph{Say it:} \emph{tolṇo}, once more
  \item \emph{Say it:} say \emph{tolṇo}
  \item \emph{Recall:} say the Marwadi for to reach, then the Marwadi for to fly, then say \emph{tolṇo} again
\end{itemize}

\begin{tcolorbox}[breakable,colback=teal!4,colframe=teal!35!black,title={Before you move on}]
What does \mw{तोलणो} mean? (\textquotedblleft{}To weigh\textquotedblright{}.) Say it once more.
\end{tcolorbox}
